\lecture{5}
\title{DT Fourier Series}

\begin{document}

\subsection{Lecture: DT Fourier Series}

Recall that for CT signals, Fourier series are \emph{infinite} sums.

\[ f(t) = \sum_{k = 0}^{\infty} \left( c_k \cos(k \omega_0 t) + d_k \sin(k \omega_0 t) \right)
= \sum_{k = -\infty}^{\infty} a_k e^{j k \omega_0t}. \]

However, for DT signals, the period $N$ is always an integer.
And when $N$ is an integer, something special happens.

\begin{align*}
    \frac{1}{a_{k + N}}\left((k + N)^{\text{th}} \text{ term}\right) =
    e^{j (k + N) \Omega_0 n} & = e^{j k \Omega_0 n} \cdot e^{j N \Omega_0 n} \\
    & = e^{j k \Omega_0 n} \cdot e^{j N (2\pi / N) n} \\
    & = e^{j k \Omega_0 n} \cdot e^{2 \pi j n}  = e^{j k \Omega_0 n}
    = \frac{1}{a_k}\left(k^{\text{th}} \text{ term}\right).
\end{align*}

Every $N$ terms of the infinite sum are ``like terms''—so our \emph{infinite} series
collapses down to a \emph{finite} series!

\begin{remark}
    The above logic is unique to DT signals
    because of the final step $e^{2\pi j n} = 1$, which only holds
    assuming $n \in \mathbb{Z}$.
\end{remark}

\begin{theorem}{DT Fourier Series}
    Given any DT signal $f: \mathbb{Z} \to \mathbb{R}$ with period $N$
    and frequency $\Omega_0 = \frac{2\pi}{N}$, its Fourier series is:
    \[ f[n] = \sum_{k = 0}^{N - 1} a_k e^{j k \Omega_0 n} ~~~~~ \text{ where } ~~~~~
    a_k = \frac{1}{N} \sum_{n = 0}^{N - 1} f[n] e^{-j k \Omega_0 n} ~ \text{ for all }
    k = 0, 1, \dots, N - 1. \]
\end{theorem}

\begin{proof}
    The same argument as earlier.  The point is that
    $\sum_{n = 0}^{N - 1} e^{j(k - \ell)\Omega_0n}$
    equals $N$ if $k \equiv \ell \pmod{N}$ and $0$ otherwise. ~ $\blacksquare$
\end{proof}

Just as before, we can also consider transformations of DT signals;
the logic is the exact same as before.

\begin{theorem}{Transforming DT Signals}
    Let $F[k]$ denote the Fourier series coefficients for $f[n]$.
    Assume period $N$.
    \begin{itemize}
        \item \textbf{Linearity:} If $f_3[n] = A f_1[n] + B f_2[n]$,
        then $F_3[k] = A F_1[k] + B F_2[k]$.
        \item \textbf{Time Shift:} If $g[n] = f[n - n_0]$,
        then $G[k] = e^{- j k \Omega_0 n_0} F[k]$.
        \item \textbf{Time Reversal:} If $g[n] = f[-n]$,
        then $G[k] = F[-k]$.
        \item \textbf{Conjugate Symmetry:} If $f[n]$ is real-valued,
        then $F[k]$ and $F[-k]$ are conjugates.
    \end{itemize}
\end{theorem}

\subsection{Lecture: Pitch Perception}

Recall that when (say) a piano plays A440 Hz,
\href{https://www.youtube.com/watch?v=3IAMpH4xF9Q}{
    there are more frequencies present
} than just plain 440 Hz.  In other words, most of the coefficients
of a piano's Fourier series are nonzero!

But the existence of, say, 880 Hz and 1320 Hz frequencies
doesn't make us hear the piano's pitch as anything other than 440 Hz\dots

\begin{quote}
    \textbf{Question.} To what extent is perceived pitch influenced by upper harmonic frequencies?
\end{quote}

In the setup below, we pass a jet of air through a spinning disk with holes.
By changing the hole placement (but keeping the disk speed constant),
we can produce different DT signals with the same fundamental frequency.

\begin{center}
    \frame{\includegraphics[width=0.45\textwidth]{lecture-05/siren}}
\end{center}

By computing Fourier coefficients for the DT signal $f_i[n]$,
we can find the strength of the ``intended'' pitch by reading off $F_i[1]$.

\begin{center}
    \frame{\includegraphics[width=0.45\textwidth]{lecture-05/siren-DT}}
\end{center}

Mathematically, the intended pitch appears to be weakest around $f_5[n]$.
In particular, when the two holes are \emph{precisely} opposite each other,
the DT signal \emph{exactly} matches $f_5[n]$,
and the perceived pitch ends up being twice the intended pitch.

The surprising part, however, is that the perceived pitch of $f_4[n]$ and $f_6[n]$
still matches the intended pitch.  Only in the exact case of $f_5[n]$
is the perceived pitch wrong.

Let's say we alter the signals by setting only the coefficient $F_i[1] \to 0$,
thereby removing the fundamental frequency and leaving behind only the harmonics. Even then,
\href{https://wtt.pauken.org/chapter-4-2/the-missing-fundamental}{
    the perceived pitch \emph{still} matches the intended pitch!
}

\subsection{Recitation: Computing DT Fourier Series}

As before, there are sometimes tricks for computing DT Fourier series.

\begin{problem}
    Determine the Fourier coefficients for
    $x_3[n] := \cos\left(\frac{2\pi}{32}n\right) + \sin\left(\frac{2\pi}{32}n\right)$.
\end{problem}

\begin{solution}
    The trick is to use Euler's formula: rewrite
    $\cos\left(\frac{2\pi}{32}n\right)$ and $\sin\left(\frac{2\pi}{32}n\right)$
    as complex exponentials.  It's not hard.
\end{solution}

\end{document}